\section{What the primitive costs}
\label{sec:tech:bench}

Correctness settles whether the construction is right. This section settles
whether it is affordable, which is the question that decides if the delivery
layer is buildable at all.

\subsection{How the numbers are produced}
\label{sec:tech:bench:harness}

Every measurement is emitted as one JSON object per line, carrying the fields
fixed by the design draft: the commit hash, the host, the network profile, the
parameter tuple, the pipeline stage and phase, wall time, and bytes sent. The
commit hash, host and timestamp are filled by the emitter rather than by the
caller, so a row cannot exist without provenance. The phase is validated
against the permitted set on the way out, because a misspelled phase produces a
row that a plotting script silently drops, and that is discovered days later.

Two conventions are worth stating. Wall time is recorded \emph{per operation},
with the batch size carried alongside so batch time is recoverable. And one row
is written per repetition, never an aggregate, because the raw file is meant to
be raw and aggregation inside the producer would hide the distribution.

The benchmark writes to standard output and its human-readable table to
standard error, with the redirection living in the build system. The binary
therefore has no file handling and no assumption about the working directory.

Guarding against the compiler optimising away the work being timed is partly
structural, since the primitive lives in a separate translation unit with
link-time optimisation switched off, so its bodies are not visible to be
elided. That would stop being true the moment anyone enabled it, so every
result is additionally accumulated into a live value that is forced to memory
and printed as a checksum. Query indices are precomputed outside the timed
region, because generating them inside it would time the generator, and a
strictly increasing sweep would give the branch predictor an unrealistically
easy task on the tree path. The whole-domain output buffer is allocated once,
outside every loop, since a freshly zeroed eight megabyte buffer per call would
measure the page fault handler rather than the primitive.

\subsection{Measured cost}
\label{sec:tech:bench:results}

\begin{table}[h]
\centering
\small
\caption{Median of five repetitions, 64-bit ring unless stated, on the
development laptop with no network involved. Read the ratios rather than the
absolute figures, for the reason given below.}
\label{tab:tech:bench}
\begin{tabular}{@{}rrrrr@{}}
\toprule
domain bits & $\Gen$ ($\mu$s) & $\Eval$ ($\mu$s) &
$\EvalFull$ ($\mu$s) & $\EvalFull$, $\bits=128$ ($\mu$s) \\
\midrule
8  & 0.40 & 0.27 & 18.4   & 52.4    \\
10 & 1.40 & 0.31 & 81.9   & 199.1   \\
11 & 1.76 & 0.44 & 281.0  & 440.8   \\
12 & 1.97 & 0.47 & 379.9  & 889.5   \\
14 & 1.74 & 0.43 & 1352.7 & 4483.8  \\
16 & 1.93 & 0.53 & 5536.0 & 14343.0 \\
\bottomrule
\end{tabular}
\end{table}

Three things follow.

\emph{The delivery layer is affordable at the scale this project targets.} At
eleven domain bits, which is the padded MovieLens-100K catalogue, one server
answering one query costs a single whole-domain evaluation, measured at roughly
$0.28$~ms. Key generation on the client is under two microseconds and the key
itself is 227 bytes. Nothing here is close to being the bottleneck.

\emph{Whole-domain evaluation earns its place.} Evaluating the entire domain
through repeated single-point evaluation is between three and four times slower
than the depth-first whole-domain pass across the range measured, because the
tree walk is shared rather than repeated from the root for every index. That
ratio is the entire justification for the second entry point existing.

\emph{Key generation is not dominated by drawing randomness.} A separate
measurement of a thirty-two byte draw from the operating system generator gives
about $0.07\,\mu$s, so the two draws a key needs account for well under a
tenth of generation time even at the smallest domain. This was worth checking
rather than assuming, because if it had gone the other way the sensible
optimisation would have been buffering randomness rather than touching the
tree, and that is not where the cost is.

\subsection{A caveat about the absolute figures}
\label{sec:tech:bench:caveat}

These numbers were taken on a laptop. Running the benchmark repeatedly in quick
succession produced absolute timings differing by as much as a factor of four
between runs, while the spread across repetitions \emph{within} a single run
stayed between $1.1$ and $1.7$. That pattern is consistent with thermal
throttling rather than with measurement noise.

The consequence is stated plainly rather than smoothed over. Ratios between
operations measured in the same run are trustworthy. Absolute figures should be
read as orders of magnitude, and the conclusions above are drawn only from
ratios and from margins wide enough that a factor of four does not change them.
Every repetition is retained in the raw file with its own timestamp, so the
drift is visible to anyone who looks rather than being averaged away. Obtaining
figures that would survive being quoted precisely needs a machine that is not
thermally constrained, and that is deferred rather than faked.
